\chapter{Literature Survey}

A review of recent research (2024--2026) was carried out across three themes directly relevant to this project: graph-based fraud detection, machine-learning-driven Anti-Money-Laundering (AML) systems, and explainable or LLM/RAG-based compliance tooling. Fourteen peer-reviewed and pre-print works were examined in total. Full bibliographic details, including author names, are provided in the References chapter; this section discusses each work by its title and contribution.

\section{Graph Neural Networks for Fraud Detection}

A substantial body of recent work applies Graph Neural Networks (GNNs) to financial fraud detection, motivated by the observation that fraud rarely occurs in isolation and is better modelled as a pattern across a network of accounts, devices, and transactions rather than as an isolated data point. \textit{``Graph Neural Networks for Financial Fraud Detection: A Review''} \cite{gnn_review2024} surveys the landscape of GNN architectures applied across fraud domains, establishing a taxonomy of node-level, edge-level, and graph-level fraud-detection tasks. Building on this foundation, \textit{``CaT-GNN: Causal Temporal Graph Neural Networks for Credit Card Fraud Detection''} \cite{catgnn2024} introduces causal and temporal reasoning into the graph representation, allowing the model to distinguish genuine causal transaction patterns from coincidental structural similarity, and reports improved performance over prior baselines on both private and public credit-card datasets. Similarly, \textit{``Detecting Credit Card Fraud via Heterogeneous Graph Neural Networks with Graph Attention''} \cite{hgnn2025} models transaction networks as heterogeneous graphs containing multiple node types -- users, merchants, and transactions -- and applies graph attention to weigh the influence of different neighbour types rather than treating every connection equally. \textit{``Harnessing Graph Neural Networks for Enhanced Fraud Detection''} \cite{gnn_enhance2025} benchmarks GNN-based approaches directly against traditional machine-learning baselines and reports consistent gains in detecting coordinated, multi-account fraud, while \textit{``TMR-GGNN: Time-Aware Multi-Relational Guided Graph Neural Network''} \cite{tmrggnn2026} further extends this line of work with a time-aware, multi-relational GNN that explicitly encodes the temporal ordering of transactions alongside multiple relationship types (e.g. shared device, shared recipient) between accounts.

\section{AML and Money-Laundering Detection}

A second cluster of literature addresses Anti-Money-Laundering detection specifically, often with an emphasis on reducing the false-positive burden that plagues real-world compliance teams. \textit{``Enhancing Anti-Money Laundering Protocols: Employing Machine Learning to Minimise False Positives and Improve Operational Cost Efficiency''} \cite{aml_ml2024}, presented at the 2024 International Conference on Advances in Artificial Intelligence, proposes a machine-learning pipeline explicitly aimed at reducing false-positive AML alerts, directly addressing the problem statement of this project. A systematic literature review, \textit{``Machine Learning Models for Money Laundering Detection in Financial Institutions''} \cite{aml_review2024}, presented at LACCEI 2024, consolidates the range of machine-learning techniques applied to money-laundering detection across the recent literature, while \textit{``Comparative Analysis of Machine Learning Algorithms for Money Laundering Detection''} \cite{aml_compare2025} provides a comparative benchmark of several classifiers (including gradient-boosted trees and neural approaches) for suspicious-transaction detection. \textit{``Leveraging Artificial Intelligence for Enhanced Anti-Money Laundering in Financial Institutions''} \cite{aml_leverage2026} reviews the practical adoption patterns of AI within bank AML programmes, highlighting the gap between what is technically demonstrated in research and what is operationally deployed. Finally, \textit{``DPxFin: Adaptive Differential Privacy for Anti-Money Laundering Detection via Reputation-Weighted Federated Learning''} \cite{dpxfin2026} addresses a complementary concern -- data privacy -- by proposing a federated-learning approach to AML detection that allows multiple institutions to collaboratively train a detection model without sharing raw customer data directly.

\section{Explainability and LLM/RAG-Based Compliance}

The third theme concerns explainability and the emerging use of large language models (LLMs) in regulatory and compliance workflows. \textit{``Secure and Transparent Banking: Explainable AI-Driven Federated Learning for Fraud Detection''} \cite{xai_fed2025} combines federated learning with SHAP- and LIME-based explainability techniques to produce fraud-detection decisions that remain both privacy-preserving and interpretable to a human reviewer. \textit{``Financial Fraud Detection Using Explainable AI and Stacking Ensemble Methods''} \cite{xai_stack2025} applies a stacking ensemble alongside SHAP explanations, demonstrating that ensemble accuracy gains need not come at the cost of interpretability. \textit{``A User-Centered Explainable AI Approach for Financial Fraud Detection''} \cite{xai_user2024} takes a user-centred approach, designing explanations specifically around the practical needs of a fraud analyst rather than around what is mathematically convenient to compute. Most relevant to the compliance-automation component of this project, \textit{``FinGuard: Detecting Financial Regulatory Non-Compliance in LLM Interactions''} \cite{finguard2026} proposes a system that uses LLMs to detect regulatory non-compliance within financial workflows, illustrating both the promise and the current limitations of applying generative AI directly to regulatory text.

\section{Limitations of the Existing Works}

While the reviewed literature demonstrates strong individual results, several consistent gaps emerge when the works are considered together, and these gaps directly motivate the design of this project.

\begin{itemize}
    \item \textbf{Detection models are studied in isolation.} The GNN-based fraud-detection papers \cite{gnn_review2024, catgnn2024, hgnn2025, gnn_enhance2025, tmrggnn2026} each optimise a single detection model and evaluate it against a fixed benchmark dataset. None of them integrate the detection model into a broader investigation or compliance workflow -- a trained model that outputs a risk score is only one part of what a real fraud or compliance team actually needs.
    \item \textbf{AML systems remain largely siloed per institution.} The AML-focused works \cite{aml_ml2024, aml_review2024, aml_compare2025, aml_leverage2026, dpxfin2026} reduce false positives and improve detection accuracy, but generally operate on a single institution's transaction data without a shared, cross-source knowledge graph that links transactions to KYC records, device fingerprints, and behavioural history in one structure.
    \item \textbf{Explainability is applied to one model at a time.} The XAI-focused papers \cite{xai_fed2025, xai_stack2025, xai_user2024} add SHAP- or LIME-based interpretability to a single classifier, but none extend this to a multi-stage decision pipeline, where an analyst may need an explanation that spans several models and a knowledge-graph traversal, not just one.
    \item \textbf{LLM/RAG-based compliance tooling is not integrated with real-time detection.} FinGuard \cite{finguard2026} demonstrates that LLMs can be applied to regulatory-compliance checking, but this and similar tools operate as a separate layer, disconnected from the underlying fraud or AML detection models that would actually trigger a compliance review.
    \item \textbf{Validation is largely restricted to public or synthetic benchmark datasets} (for example, the widely-used Kaggle credit-card fraud dataset), rather than at the scale and complexity of a real bank's multi-source data. This is an acknowledged limitation of the present project as well, and is discussed further in Chapter~6.
\end{itemize}

Taken together, no reviewed work combines a unified, cross-source knowledge graph, a real-time and explainable risk-scoring model, and a RAG-grounded compliance-automation pipeline within a single deployable system. This combination is precisely the gap that the proposed platform, described in Chapter~5, is designed to address.
